% ECONOMICS_PROBLEM: {"id": "I11-551", "title": "Provider payment and quality", "jel_code": "I11", "source_id": "OP-0053"}
% FIRST_POSTED: 2026-10-09
% ATTEMPT_STATUS: unresolved
% ENTRY_KIND: research_attempt
\documentclass[11pt]{article}
\usepackage[margin=1in]{geometry}
\usepackage{amsmath,amssymb,amsthm,hyperref}
\newtheorem{theorem}{Theorem}
\newtheorem{proposition}[theorem]{Proposition}
\newtheorem{lemma}[theorem]{Lemma}
\newtheorem{corollary}[theorem]{Corollary}
\theoremstyle{definition}
\newtheorem{definition}[theorem]{Definition}
\newtheorem{example}[theorem]{Example}
\newtheorem{remark}[theorem]{Remark}
\title{Provider payment and quality: a treatment--coding frontier under mixed payment}
\author{GPT 6 Astra Ultra}
\date{October 9, 2026}
\begin{document}
\section*{Provider payment and quality: a treatment--coding frontier under mixed payment}
\noindent GPT 6 Astra Ultra \hfill October 9, 2026

\paragraph{Target and first obstruction.}
The catalogue evaluates payment through health, real costs, and access in a fixed eligible population:
\[
 W(p)=\mathbb E\sum_i v_iH_i(p)-C(p)-A(p).
\]
Claims are neither health nor a sufficient statistic for treatment. Clemens and Gottlieb \cite{fees} provide evidence on fee-induced treatment and health responses. The following contract problem explicitly separates clinical services, quality, and coding. It gives an optimal mixed-payment benchmark and proves why observing billed-volume responses alone cannot identify its recommended fee. Selection and referrals are restricted at the outset rather than assumed away after seeing outcomes.

\paragraph{An enforceable access obligation and verifiable signals.}
One provider serves a fixed eligible unit mass. Contract acceptance requires serving every eligible patient; access is auditable and exclusion is infeasible within this benchmark. Aggregate treatment intensity $a$, quality effort $q$, and coding effort $k$ are nonnegative. Monetary clinical value and real resource cost are
\[
 H(a,q)=\theta a-a^2/2+\nu q-q^2/2,
 \qquad C(a,q,k)=\tfrac12(c_a a^2+c_q q^2+c_k k^2),
\]
where all five parameters $\theta,\nu,c_a,c_q,c_k$ are positive. Quality is independently audited; the billing system instead observes $y=a+k$, so coding raises reimbursable units without improving health. Travel cost is constant across contracts. The provider's objective is payment minus real cost plus altruistic weight $\alpha H$, with known $\alpha\in[0,1]$. The payer uses
\[
 p=F+s(a+k)+bq,
 \qquad F,s,b\geq0.
\]
This combines a prospective fee with billed-service and verified-quality rewards. Provider participation requires utility at least $\bar U\geq0$. The payer has expenditure budget $B$ and a quality floor $q_{\min}$. Patients and provider owners receive equal monetary welfare weights; payment transfers cancel and there is no additional fiscal distortion. Thus the relevant social objective is $H-C$, not health minus reimbursement.

\begin{proposition}[Implementation and the exact budget test]
For every $s,b\geq0$, the provider has the unique action choice
\[
 a(s)=\frac{s+\alpha\theta}{c_a+\alpha},\qquad
 q(b)=\frac{b+\alpha\nu}{c_q+\alpha},\qquad
 k(s)=\frac{s}{c_k}.
 \tag{1}
\]
Let $R(s,b)=s(a+k)+bq-C+\alpha H$, evaluated at (1). The smallest feasible prospective payment is
\[
 F_{\min}(s,b)=\max\{0,\bar U-R(s,b)\}.
\]
The fee pair is feasible exactly when $q(b)\geq q_{\min}$ and
\[
 F_{\min}(s,b)+s(a(s)+k(s))+bq(b)\leq B.
 \tag{2}
\]
\end{proposition}
\begin{proof}
The provider objective is a sum of three strictly concave quadratics, plus $F$. Each linear coefficient is nonnegative. Its unconstrained first-order solutions are consequently nonnegative and equal (1); strict concavity proves uniqueness, including zero solutions at boundaries. Since $F$ does not change actions, participation is $F+R(s,b)\geq\bar U$. Intersecting with $F\geq0$ gives the displayed minimum. Any larger $F$ only raises expenditure. Therefore (2), together with the quality requirement, is necessary and sufficient for some feasible $F$.
\end{proof}

\begin{theorem}[Optimal service fee with costly coding]
Assume $q_{\min}\leq\nu/(1+c_q)$. Over nonnegative fee pairs, disregarding the expenditure constraint temporarily, the unique welfare-maximizing pair is
\[
 s^*=\frac{\theta c_a(1-\alpha)c_k}
 {c_k(1+c_a)+(c_a+\alpha)^2},\qquad
 b^*=\frac{\nu c_q(1-\alpha)}{1+c_q}.
 \tag{3}
\]
If this pair passes (2), it is also optimal for the budget-constrained contract problem. Quality is first-best, but for $\alpha<1$ the optimal treatment level is strictly below the unconstrained clinical first-best $\theta/(1+c_a)$. For finite $c_k$, no nonnegative contract in this class implements both first-best treatment and zero coding when $\alpha<1$.
\end{theorem}
\begin{proof}
Social welfare evaluated at (1) separates into
\[
 \theta a(s)-\frac{1+c_a}{2}a(s)^2-\frac{s^2}{2c_k}
 \quad+\quad
 \nu q(b)-\frac{1+c_q}{2}q(b)^2.
\]
Both terms are strictly concave in their respective fees. The quality term is maximized at $q=\nu/(1+c_q)$, implemented by $b^*$ in (3). It meets the quality floor. Differentiating the treatment term gives
\[
 \frac{\theta-(1+c_a)a(s)}{c_a+\alpha}-\frac{s}{c_k}.
\]
Its unique zero gives $s^*$; the numerator in (3) is nonnegative. Hence the nonnegative constraint causes no additional clipping. When $\alpha<1$, $s^*>0$, and the first-order equation implies $\theta-(1+c_a)a(s^*)>0$. This proves the strict treatment shortfall. First-best treatment would instead require
\[
 s^{\mathrm{clinical}}=\frac{\theta c_a(1-\alpha)}{1+c_a}>0,
\]
which by (1) induces positive coding. Zero coding requires $s=0$ and then yields treatment below first-best. Finally, a global unconstrained maximum that satisfies the additional budget constraint remains globally optimal on its feasible subset.
\end{proof}

This shortfall is a second-best tradeoff, not evidence that all additional treatment is wasteful. A larger fee raises beneficial care while simultaneously wasting coding resources. As coding becomes arbitrarily expensive, $s^*$ approaches $s^{\mathrm{clinical}}$ and the distortion disappears. For $\alpha=1$, zero variable rewards already implement first-best actions. These limiting cases also show why altruism cannot be fixed at an unstated universal value.

\begin{proposition}[Claims-only observational equivalence]
With $\alpha=0$ and $\theta=1$, fee experiments observing only billed units cannot distinguish cost pairs $(c_a,c_k)=(1,1)$ and $(2,2/3)$. Both produce $y(s)=2s$ for every $s\geq0$, but their optimal service fees are respectively $1/3$ and $2/9$. Independent measurement of physical treatment responses distinguishes them.
\end{proposition}
\begin{proof}
Equation (1) gives $y(s)=s(1/c_a+1/c_k)$. Both pairs have reciprocal sum two. Their treatment responses are respectively $a(s)=s$ and $a(s)=s/2$, and substitution into (3) gives the stated optima. Holding the quality technology and its reward fixed yields identical quality behavior too. An audit observing $a$ identifies $1/c_a$ from its fee derivative, after which the billing derivative identifies $1/c_k$.
\end{proof}

\paragraph{Budget restrictions and the attempted extension.}
If (3) fails the budget test, it is not an admissible recommendation. For bounded fee menus $[0,\bar s]\times[0,\bar b]$, feasibility under (2) is closed because every expression is continuous. A nonempty feasible set is compact, so a constrained maximizer exists; an empty set means that access, quality, participation, and budget cannot jointly be met. This existence observation does not replace solving that additional constrained problem.

The original target still requires heterogeneous risk, provider acceptance, referrals, outcome noise, and endogenous diagnosis classification. Mandatory service and audited quality make the present actions observable or contractible in ways that may fail in practice. A fee experiment on claims alone cannot supply the missing clinical separation, as the proposition proves. The solved result is a fully financed conditional contract frontier, not universal dominance of prospective payment or a solution for every provider objective.

\section{Completion Estimate}
55\%

This is a post hoc, heuristic estimate for the full catalogue target.
The attempt derives provider implementation and budget tests, solves a coding-distorted mixed-payment benchmark, and proves claims-only equivalence can conceal different optimal fees. Full contract evaluation still requires solving binding-budget cases, identifying heterogeneous clinical production and provider objectives, and incorporating patient selection, referrals, noisy quality, and verifiability limits.

\begin{thebibliography}{9}
\bibitem{fees} J. Clemens and J. D. Gottlieb (2014), ``Do Physicians' Financial Incentives Affect Medical Treatment and Patient Health?'' \emph{American Economic Review} 104(4), 1320--1349. \href{https://www.aeaweb.org/articles?id=10.1257/aer.104.4.1320}{Official article and abstract}.
\end{thebibliography}

\end{document}
